\chapter{Tao Yuanming and Fields-and-Gardens Poetry}

\section{Background: the Six Dynasties}

After the Han collapsed (220), China was divided for more than three and a half
centuries. In 317 the Jin court fled south to Jiankang (modern Nanjing), and the north was
ruled by non-Han dynasties. In the south, an aristocratic culture flourished: philosophy of
``Dark Learning'' (\zht{玄學}, Neo-Daoism), Buddhism, ``pure conversation'' (\zht{清談}),
calligraphy, and poetry that grew increasingly ornate.

Two genres of nature poetry emerge in this period:
\begin{itemize}
  \item \textbf{Landscape poetry} (\zht{山水詩} \emph{shānshuǐ shī}, ``mountains and
        waters''), founded by Xie Lingyun \zht{謝靈運} (385--433): wild scenery, travel,
        elaborate parallel description, Buddhist/Daoist insight.
  \item \textbf{Fields-and-gardens poetry} (\zht{田園詩} \emph{tiányuán shī}), founded by
        Tao Yuanming: farming, the village, home, simplicity.
\end{itemize}

\section{Tao Yuanming \zht{陶淵明} (365?--427)}

Also known as Tao Qian \zht{陶潛}. From a family of declining minor gentry, he took a
series of minor official posts to support his family. In 405, after about eighty days as
magistrate of Pengze, he resigned, famously refusing to ``bow and scrape for five pecks of
rice'' (the official salary) before a visiting inspector. He spent the rest of his life
farming near Mount Lu, poor, drinking, writing.

\begin{important}[Tao Yuanming in one paragraph]
Tao Yuanming is the model of the \textbf{recluse} (\zht{隱士}) who chooses integrity
and freedom over office. His style is \textbf{plain} (\zht{平淡} \emph{píngdàn}, ``even and
bland''): everyday words, few allusions, little ornament, at a time when poetry valued
ornament. Underappreciated in his own time (Zhong Rong's \emph{Grades of Poets} ranks him
only middle grade), he was elevated by Tang and especially Song poets (Su Shi wrote poems
``matching'' all of Tao's) into one of the greatest Chinese poets.
\end{important}

\begin{center}
\begin{tikzpicture}[box/.style={draw, thick, rounded corners, align=center, text width=3.6cm,
  minimum height=1.6cm, font=\small}]
  \node[box, fill=gray!15] (off) at (0,0) {\textbf{Office} \zht{仕}\\ the ``dusty net''\\ cage, bowing, salary};
  \node[box, fill=green!15] (rec) at (7,0) {\textbf{Reclusion} \zht{隱}\\ fields and gardens\\ the ``natural'' \zht{自然}};
  \draw[-{Stealth}, very thick] (off) -- node[above, font=\scriptsize] {405: ``Return home!''}
    node[below, font=\scriptsize] {\zht{歸去來兮}} (rec);
  \node[font=\scriptsize, align=center] at (0,-1.6) {caged bird, fish in a pond};
  \node[font=\scriptsize, align=center] at (7,-1.6) {bird returning to the old wood};
\end{tikzpicture}
\end{center}

\begin{aside}
The Confucian ideal was to serve when the world is well governed and withdraw when it is
not. Reclusion was therefore a respectable \emph{political} statement, not merely a taste
for the country. Tao's farming was real, though: his poems mention hunger, failed harvests
and a house fire.
\end{aside}

\section{``Returning to Live in the Fields'' \zht{歸園田居} (no.\ 1 of 5)}

\begin{center}
\begin{tabular}{p{0.3\textwidth}p{0.6\textwidth}}
\toprule
\zht{少無適俗韻，} & When young I had no tune that suited the common world; \\
\zht{性本愛丘山。} & by nature I loved the hills and mountains. \\
\zht{誤落塵網中，} & By mistake I fell into the dusty net \\
\zht{一去三十年。} & and was gone for thirty years. \\
\zht{羈鳥戀舊林，} & The caged bird longs for its old wood; \\
\zht{池魚思故淵。} & the fish in the pond thinks of its former depths. \\
\zht{開荒南野際，} & I clear waste land at the edge of the southern wilds, \\
\zht{守拙歸園田。} & keeping to my clumsiness, returned to fields and garden. \\
\zht{方宅十餘畝，} & Ten-odd \emph{mu} round my house, \\
\zht{草屋八九間。} & a thatched cottage of eight or nine rooms. \\
\zht{榆柳蔭後簷，} & Elms and willows shade the back eaves; \\
\zht{桃李羅堂前。} & peach and plum stand in rows before the hall. \\
\zht{曖曖遠人村，} & Dim, dim, the villages of distant people; \\
\zht{依依墟里煙。} & soft and lingering, the smoke above the hamlets. \\
\zht{狗吠深巷中，} & A dog barks somewhere deep in the lanes; \\
\zht{雞鳴桑樹顛。} & a cock crows at the top of a mulberry tree. \\
\zht{戶庭無塵雜，} & My door and yard are free of dust and clutter; \\
\zht{虛室有餘閒。} & in the empty rooms there is leisure to spare. \\
\zht{久在樊籠裏，} & Long I was inside a cage; \\
\zht{復得返自然。} & now again I can return to what is so of itself. \\
\bottomrule
\end{tabular}
\end{center}

\subsection{Analysis}
\begin{itemize}
  \item \textbf{Structure.} Lines 1--8: the frame (mistake, cage, return). Lines 9--16: an
        inventory of the homestead, moving outward from house to trees to distant villages
        and back to sounds. Lines 17--20: the conclusion returns to the cage image.
  \item \textbf{Parallelism} is everywhere but feels unforced: caged bird / pond fish;
        elms-willows / peach-plum; dog barking / cock crowing.
  \item \textbf{``Clumsiness''} (\zht{拙}): a key Tao word. He is unfit for the clever
        world, and proud of it.
  \item \textbf{Allusion, lightly worn.} ``Dog barks deep in the lanes, cock crows atop
        the mulberry'' echoes a Han \emph{yuefu} line, and the image of neighbouring
        villages hearing each other's dogs and chickens comes from the \emph{Laozi}'s
        ideal small community.
  \item \textbf{\zht{自然} \emph{zìrán}} ``self-so, natural''. In modern Chinese it means
        ``nature''; here it is the Daoist ideal of spontaneity, being what one is.
  \item \textbf{``Thirty years''}: he actually served about thirteen, and some scholars
        emend to \zht{十三年} ``thirteen years''. A useful reminder that texts were transmitted by hand.
\end{itemize}

\section{``Drinking Wine'' \zht{飲酒} (no.\ 5 of 20)}

Tao's most famous poem.

\begin{center}
\begin{tabular}{p{0.24\textwidth}p{0.26\textwidth}p{0.4\textwidth}}
\toprule
Original & Pinyin & My translation \\
\midrule
\zht{結廬在人境，} & jié lú zài rén jìng & I built my hut among other people, \\
\zht{而無車馬喧。} & ér wú chē mǎ xuān & yet hear no noise of carriages and horses. \\
\zht{問君何能爾？} & wèn jūn hé néng ěr & You ask me how this can be? \\
\zht{心遠地自偏。} & xīn yuǎn dì zì piān & The mind far off, the place grows remote of itself. \\
\zht{采菊東籬下，} & cǎi jú dōng lí xià & Picking chrysanthemums by the eastern hedge, \\
\zht{悠然見南山。} & yōurán jiàn nán shān & far and calm, I see the southern hills. \\
\zht{山氣日夕佳，} & shān qì rì xī jiā & The mountain air is lovely at sunset; \\
\zht{飛鳥相與還。} & fēi niǎo xiāng yǔ huán & birds in flight return together. \\
\zht{此中有真意，} & cǐ zhōng yǒu zhēn yì & In this there is a true meaning --- \\
\zht{欲辨已忘言。} & yù biàn yǐ wàng yán & I would explain it, but have already forgotten the words. \\
\bottomrule
\end{tabular}
\end{center}

\subsection{Analysis}
\begin{itemize}
  \item \textbf{Reclusion of the mind.} He does not live on a remote mountain but ``among
        people''. Reclusion is a state of mind (\zht{心遠}), not a location: an idea that
        made Tao a model for officials who could not afford to resign.
  \item \textbf{\zht{見} \emph{jiàn} ``see''.} Su Shi argued that the variant \zht{望}
        \emph{wàng} ``gaze at'' ruins the line: with \zht{見} the mountain simply comes
        into view, unsought, while he is busy picking flowers. The difference between
        seeing and looking is the whole poem.
  \item \textbf{Chrysanthemum}: flower of autumn and longevity; after this poem, forever
        Tao's flower.
  \item \textbf{Birds returning} (\zht{還}) at dusk mirror the poet's own return home.
  \item \textbf{Ending.} An allusion to \emph{Zhuangzi}: ``words exist for meaning; once you
        have the meaning you can forget the words.'' The poem ends by declining to explain
        itself.
\end{itemize}

\begin{important}[Exam ID: \zht{采菊東籬下，悠然見南山}]
Tao Yuanming, ``Drinking Wine'' no.~5, Eastern Jin. Fields-and-gardens poetry; the
recluse's spontaneous, unforced harmony with nature; \zht{見} vs \zht{望}.
\end{important}

\section{Other works}

\begin{itemize}
  \item \textbf{``Return Home!''} (\zht{歸去來兮辭}, 405): a rhapsody in \emph{Chuci}
        style (with \zht{兮}) celebrating his resignation.
  \item \textbf{``Peach Blossom Spring''} (\zht{桃花源記}): a prose story with an
        attached poem. A fisherman follows a stream of peach blossoms to a hidden village
        whose people fled the wars of the Qin and have lived in peace for centuries,
        unaware of the Han and Wei. When he leaves and tries to find it again, he cannot.
        \zht{桃花源} became the Chinese word for utopia.
  \item \textbf{Self-elegies}: Tao wrote ``Elegies for Myself'' (\zht{擬挽歌辭}),
        contemplating his own funeral with calm humour.
\end{itemize}

\section*{Summary}
\begin{itemize}
  \item Six Dynasties: landscape poetry (Xie Lingyun) and fields-and-gardens poetry (Tao
        Yuanming).
  \item Tao Yuanming resigned office in 405; became the archetypal recluse.
  \item Plain style, light allusion, Daoist \zht{自然}; caged bird / returning bird.
  \item ``Drinking Wine'' no.~5: reclusion of the mind; chrysanthemums; words forgotten.
  \item Undervalued in his lifetime, canonised by Tang and Song poets.
\end{itemize}
